\chapter{Concurrency Control}
\label{ch:concurrency}

Isolation levels say \emph{what} must hold; concurrency control is \emph{how}.
This chapter covers locking, the approach OntoDB takes: two-phase locking,
multi-granularity locks, deadlock handling, phantoms, and the separate problem
of protecting the B+ tree's own structure.

\section{Two-phase locking}

A transaction takes a \term{shared} (S) lock to read an item and an
\term{exclusive} (X) lock to write it. S is compatible with S; X is compatible
with nothing. \term{Two-phase locking} (2PL) adds one rule: a transaction may
not acquire any lock after it has released one. Its life has a growing phase
and a shrinking phase.

2PL guarantees conflict-serializability: the order in which transactions reach
their \term{lock point} (the moment they hold all their locks) is an equivalent
serial order.

Plain 2PL has a flaw. If $T_1$ releases an X lock before committing and $T_2$
reads the value, then $T_1$ aborts, $T_2$ has read a value that never existed,
and must abort too: a \term{cascading abort}. \term{Strict 2PL} holds every
X lock until commit or abort, which rules this out.

\begin{keyidea}
Exclusive locks are held to the end of the transaction at every isolation
level, because releasing one early lets others see a write that may still be
undone. The isolation levels differ only in how long \emph{shared} locks are
held, and whether ranges are locked.
\end{keyidea}

\begin{table}[ht]
\centering
\begin{tabular}{@{}ll@{}}
\toprule
Level & Shared locks \\
\midrule
Read uncommitted & none \\
Read committed & released at the end of each statement \\
Repeatable read & held until commit \\
Serializable & held until commit, plus range or predicate locks against phantoms \\
\bottomrule
\end{tabular}
\caption{Isolation levels as locking rules, given strict 2PL for writes.}
\end{table}

\section{Locking at several granularities}

Locking every triple is precise but expensive for a query that reads the whole
store; locking the whole store is cheap but serializes everything. Gray et al.\
\cite{gray1976} let a transaction choose its granularity, using
\term{intention} locks on the coarser levels: before locking a key in S, take
IS (``intends to share something below'') on the index and on the database.

\begin{table}[ht]
\centering
\begin{tabular}{@{}lccccc@{}}
\toprule
Held \textbackslash{} Requested & IS & IX & S & SIX & X \\
\midrule
IS  & yes & yes & yes & yes & --- \\
IX  & yes & yes & --- & --- & --- \\
S   & yes & --- & yes & --- & --- \\
SIX & yes & --- & --- & --- & --- \\
X   & --- & --- & --- & --- & --- \\
\bottomrule
\end{tabular}
\caption{Multi-granularity lock compatibility. SIX is S on the whole node plus the intention to write below it.}
\label{tab:compat}
\end{table}

The point of intention locks: a transaction that wants to lock a whole index in
S can check one lock on the index, instead of every key below it, because any
writer below must already hold IX there.

\section{Deadlock}

$T_1$ holds a lock $T_2$ wants, and $T_2$ holds a lock $T_1$ wants: both wait
forever. Engines either \emph{prevent} deadlock (order locks, or abort on
conflicts by timestamp, as in wait--die and wound--wait) or \emph{detect} it:
maintain a \term{waits-for graph} with an edge $T_i \to T_j$ when $T_i$ waits for
a lock $T_j$ holds, find a cycle, and abort one transaction in it, the
\term{victim}. Choosing the youngest transaction (the one that has done the least
work) is a common and deterministic policy.

\section{Phantoms}

Repeatable read locks the triples a transaction read. A phantom is a triple it
did \emph{not} read, because it did not exist yet; no lock covers it. Preventing
phantoms needs a lock on what \emph{could} match, not on what does:

\begin{description}[leftmargin=1.2em]
\item[Predicate locks] lock the query's condition itself. Precise, but testing whether two predicates overlap is expensive in general \cite{eswaran}.
\item[Key-range (next-key) locks] lock the gap between index keys, so an insert into the scanned range must wait. Practical, and natural on a B+ tree.
\end{description}

\section{Latches are not locks}

A B+ tree split rewrites several pages at once. Even a single-transaction
engine with many threads must stop two threads from splitting the same page at
the same time. That protection is a \term{latch}: a short-term reader--writer
mutex on a page, held for the duration of one tree operation, never for a
transaction. Locks protect \emph{logical contents} for transactional isolation;
latches protect \emph{physical structure} for thread safety.

The classic latch protocol for B+ trees is \term{crabbing}: latch the child,
and release the parent if the child is \emph{safe}, meaning it cannot split (for
an insert) or merge (for a delete) as a result of this operation. A thread
therefore holds latches only on the part of the path a structure change could
reach. An \term{optimistic} variant takes shared latches down to the leaf,
latches only the leaf exclusively, and restarts pessimistically if the leaf
turns out to need a split. Lehman and Yao's B-link tree \cite{lehmanyao} goes
further and lets readers run without latching the path at all, by following
right-links when a split has moved their key.

\section{The alternative: multiversion concurrency}

Locking makes readers and writers wait for each other. \term{Multiversion
concurrency control} (MVCC) keeps old versions of items, so a reader sees a
consistent snapshot as of its start time and never waits for a writer. Writers
still conflict with writers. The cost is storing old versions and cleaning them
up once no reader can need them. Most modern systems use MVCC with snapshot
isolation by default.

\begin{ontodb}[: concurrency control]
OntoDB uses strict 2PL with five lock modes on three granularities: the
database, an index (\code{"spo"}, \code{"pos"}, \code{"osp"}), and a key, with
FIFO grant order and in-place upgrades that do not jump ahead of waiters. A
blocked shell session prints why it is waiting, for example
\code{session 2 waiting: X lock held by txn 5}. A background thread builds the
waits-for graph and aborts the \emph{youngest} transaction in each cycle (the
larger id), so deadlock outcomes are deterministic and testable. Serializable
adds one phantom mechanism, chosen and documented in the code. B+ tree
operations use latch crabbing, and the whole feature runs under
ThreadSanitizer. MVCC is a deliberate future comparison, not part of the plan.
\end{ontodb}

\begin{status}
Concurrency control is plans 7.1 to 7.6: crabbing, transactions, the lock
manager, isolation, deadlock detection, and phantoms. The isolation harness
already exists and checks itself.
\end{status}

\section*{Further reading}
Gray et al.\ \cite{gray1976}; Eswaran et al.\ \cite{eswaran}; Lehman and Yao
\cite{lehmanyao}; Bayer and Schkolnick \cite{bayerschkolnick}; Petrov
\cite{petrov}, chapters 5 and 6.
